% ============================================================================
% DEPENDENCIES
% ============================================================================

\section{Dependency Relationships}
\label{sec:dependencies}

This section analyzes the dependency relationships between modules and files in the \ctx{} codebase.

\subsection{File-Level Dependencies}

The following diagram shows the primary dependency relationships between source files, based on function calls and imports.

\begin{figure}[h]
\centering
\begin{tikzpicture}[
    node distance=1.5cm,
    file/.style={rectangle, draw, rounded corners, minimum width=2.5cm, minimum height=0.7cm, align=center, fill=blue!10, font=\small},
    core/.style={file, fill=orange!20},
    parser/.style={file, fill=green!20},
    db/.style={file, fill=purple!20},
    arrow/.style={-{Stealth[length=2mm]}, thick, color=gray!70},
    label/.style={font=\tiny, color=gray}
]

% Core files
\node[core] (main) {\file{main.rs}};
\node[core, below left=1.5cm and 1cm of main] (cli) {\file{cli.rs}};

% Database layer
\node[db, below=2cm of main] (schema) {\file{db/schema.rs}};
\node[db, left=0.5cm of schema] (models) {\file{db/models.rs}};

% Parser layer
\node[parser, right=3cm of schema] (parsermod) {\file{parser/mod.rs}};
\node[parser, above right=0.5cm and 0.3cm of parsermod] (rust) {\file{rust.rs}};
\node[parser, right=0.3cm of parsermod] (ts) {\file{typescript.rs}};
\node[parser, below right=0.5cm and 0.3cm of parsermod] (py) {\file{python.rs}};
\node[parser, below=0.5cm of py] (sol) {\file{solidity.rs}};

% Index and analytics
\node[file, below left=1cm and 0cm of schema] (index) {\file{index/mod.rs}};
\node[file, below right=1cm and 0cm of schema] (analytics) {\file{analytics/mod.rs}};

% Context generation
\node[file, left=3cm of schema] (walker) {\file{walker.rs}};
\node[file, below=0.5cm of walker] (output) {\file{output.rs}};
\node[file, below=0.5cm of output] (formatter) {\file{formatter.rs}};
\node[file, left=0.5cm of formatter] (tree) {\file{tree.rs}};

% Embeddings
\node[file, fill=yellow!20, below=2.5cm of analytics] (embed) {\file{embeddings/mod.rs}};
\node[file, fill=yellow!20, left=0.3cm of embed] (local) {\file{local.rs}};
\node[file, fill=yellow!20, right=0.3cm of embed] (openai) {\file{openai.rs}};

% Arrows from main
\draw[arrow] (main) -- node[label, above, sloped] {8} (schema);
\draw[arrow] (main) -- node[label, above, sloped] {8} (parsermod);
\draw[arrow] (main) -- (cli);
\draw[arrow] (main) -- node[label, above, sloped] {7} (index);
\draw[arrow] (main) -- node[label, above, sloped] {5} (analytics);
\draw[arrow] (main) -- (walker);
\draw[arrow] (main) -- (output);
\draw[arrow] (main) -- (embed);

% Index dependencies
\draw[arrow] (index) -- node[label, above, sloped] {11} (parsermod);
\draw[arrow] (index) -- node[label, left] {6} (schema);
\draw[arrow] (index) -- (walker);

% Parser dependencies
\draw[arrow] (rust) -- (parsermod);
\draw[arrow] (ts) -- node[label, above] {13} (parsermod);
\draw[arrow] (py) -- (parsermod);
\draw[arrow] (sol) -- (parsermod);

% Output dependencies
\draw[arrow] (output) -- node[label, left] {8} (formatter);
\draw[arrow] (output) -- (tree);

% Analytics dependencies
\draw[arrow] (analytics) -- (schema);

% Embed dependencies
\draw[arrow] (local) -- (embed);
\draw[arrow] (openai) -- (embed);

\end{tikzpicture}
\caption{File-level dependency graph (numbers indicate call count)}
\label{fig:dependencies}
\end{figure}

\subsection{Key Dependency Chains}

\subsubsection{Main Entry Point}

The \file{main.rs} file serves as the central hub, with dependencies on most major modules:

\begin{verbatim}
main.rs ──┬──► db/schema.rs (8 calls)
          ├──► parser/mod.rs (8 calls)
          ├──► index/mod.rs (7 calls)
          ├──► analytics/mod.rs (5 calls)
          ├──► walker.rs
          ├──► output.rs
          ├──► formatter.rs
          └──► embeddings/mod.rs
\end{verbatim}

\subsubsection{Indexing Pipeline}

The indexing pipeline flows from file discovery through parsing to database storage:

\begin{verbatim}
index/mod.rs ──┬──► walker.rs (file discovery)
               ├──► parser/mod.rs (symbol extraction)
               │    └──► language-specific parsers
               └──► db/schema.rs (persistence)
\end{verbatim}

\subsubsection{Context Generation Pipeline}

Context generation is simpler, flowing from file selection through formatting:

\begin{verbatim}
output.rs ──┬──► walker.rs (file discovery)
            ├──► tree.rs (tree visualization)
            └──► formatter.rs (output formatting)
\end{verbatim}

\subsection{Parser Interdependencies}

An interesting pattern emerges in the parser modules---they have bidirectional dependencies:

\begin{verbatim}
typescript.rs ◄──► python.rs ◄──► rust.rs ◄──► solidity.rs
\end{verbatim}

This pattern likely arises from:
\begin{enumerate}
    \item Shared test fixtures imported across parsers
    \item Common type definitions used in all parsers
    \item Potential for further abstraction into a shared parser trait
\end{enumerate}

\subsection{Module Coupling Analysis}

\begin{table}[h]
\centering
\begin{tabular}{lrrr}
\toprule
\textbf{Module} & \textbf{Fan-Out} & \textbf{Fan-In} & \textbf{Coupling Score} \\
\midrule
\file{main.rs} & 43 & 0 & High (entry point) \\
\file{db/schema.rs} & 12 & 18 & High (central storage) \\
\file{parser/mod.rs} & 8 & 24 & High (shared utilities) \\
\file{index/mod.rs} & 18 & 8 & Medium \\
\file{formatter.rs} & 4 & 12 & Medium \\
\file{walker.rs} & 6 & 8 & Medium \\
\file{analytics/mod.rs} & 8 & 6 & Low \\
\file{embeddings/mod.rs} & 4 & 6 & Low \\
\bottomrule
\end{tabular}
\caption{Module coupling analysis}
\label{tab:coupling}
\end{table}

\subsection{Dependency Observations}

\subsubsection{Strengths}

\begin{itemize}
    \item \textbf{Clear layering} --- The codebase follows a clear layered architecture with minimal circular dependencies
    \item \textbf{Single responsibility} --- Each module has a focused purpose
    \item \textbf{Shared utilities} --- Common functionality is properly extracted to \file{parser/mod.rs}
    \item \textbf{Loose coupling for embeddings} --- The embedding system is well-isolated with a clean provider trait
\end{itemize}

\subsubsection{Areas for Improvement}

\begin{itemize}
    \item \textbf{main.rs coupling} --- The main file has high coupling to many modules; consider extracting command handlers to separate modules
    \item \textbf{Parser interdependencies} --- The bidirectional dependencies between parsers could be reduced by extracting shared test utilities
    \item \textbf{db/schema.rs size} --- With 74 symbols, this module handles many responsibilities; consider splitting into smaller focused modules
\end{itemize}

\subsection{Recommended Refactoring}

\subsubsection{Extract Command Handlers}

Move command-specific logic from \file{main.rs} to dedicated modules:

\begin{verbatim}
src/
├── commands/
│   ├── mod.rs
│   ├── index.rs      # Index command handler
│   ├── query.rs      # Query command handlers
│   ├── search.rs     # Search command handler
│   ├── embed.rs      # Embedding command handler
│   └── context.rs    # Context generation handler
├── main.rs           # Entry point only
└── ...
\end{verbatim}

\subsubsection{Split Database Module}

Separate concerns within the database layer:

\begin{verbatim}
src/db/
├── mod.rs            # Re-exports
├── schema.rs         # Schema definitions
├── symbols.rs        # Symbol CRUD operations
├── edges.rs          # Edge CRUD operations
├── search.rs         # FTS5 search operations
├── embeddings.rs     # Embedding storage
└── migrations.rs     # Schema migrations
\end{verbatim}

\subsubsection{Create Parser Trait}

Unify parser implementations behind a common trait:

\begin{lstlisting}[style=rustcode, caption={Proposed parser trait}]
pub trait LanguageParser {
    fn language(&self) -> Language;
    fn symbols_query(&self) -> &str;
    fn calls_query(&self) -> &str;
    
    fn extract_symbols(&self, tree: &Tree, source: &str) -> Vec<Symbol>;
    fn extract_edges(&self, tree: &Tree, source: &str) -> Vec<Edge>;
}

// Then in mod.rs:
pub fn get_parser(language: Language) -> Box<dyn LanguageParser> {
    match language {
        Language::Rust => Box::new(RustParser::new()),
        Language::TypeScript => Box::new(TypeScriptParser::new()),
        // ...
    }
}
\end{lstlisting}
